\section{Résultats}

\begin{tcolorbox}[colback=black!5,colframe=black,title=]
    Cette section présente les premiers résultats obtenus avec le modèle proposé. L'objectif est d'évaluer la capacité de l'algorithme à apprendre une politique de navigation efficace à partir du curriculum d'entraînement défini précédemment, puis d'analyser les comportements développés par l'agent dans les différents scénarios. Les performances sont étudiées à l'aide des critères d'évaluation présentés précédemment.
\end{tcolorbox}

%%%%%%%%%%%%%%%%%%%%%%%%%%%%%%%%%%%%%%%%%%%%%%%%%%%%%%%%%%%%%%%%%%%%%%%%%%%%%%%%%%%%%%%
%%%%%%%%%%%%%%%%%%%%%%%%%%%%%%%%%%%%%%%%%%%%%%%%%%%%%%%%%%%%%%%%%%%%%%%%%%%%%%%%%%%%%%%
%%%%%%%%%%%%%%%%%%%%%%%%%%%%%%%%%%%%%%%%%%%%%%%%%%%%%%%%%%%%%%%%%%%%%%%%%%%%%%%%%%%%%%%


\subsection{Validation de l'apprentissage}

\noindent Les courbes d'apprentissage permettent d'analyser l'évolution de la politique au cours de l'entraînement (Figure \ref{fig:valid_exp3} et Figure \ref{fig:valid_exp7}). Plus précisément, le taux de succès, le taux de collision et le taux de dépassement du temps maximal (timeout) renseignent sur la capacité de l'agent à atteindre sa cible tout en évitant les obstacles. Le temps moyen de trajet permet d'évaluer l'efficacité des trajectoires produites, tandis que la décomposition de la fonction de récompense met en évidence la contribution de chacun des termes de récompense au cours de l'apprentissage.

\vspace{10pt}

\noindent \textbf{Les premiers résultats montrent que, pour les scénarios les plus simples (expériences 1 à 4), le taux de réussite augmente progressivement jusqu'à atteindre des valeurs proches de 100\%, tandis que les taux de collision et de timeout diminuent rapidement.} L'évolution du temps moyen de trajet met également en évidence les différentes phases de l'apprentissage. Au début de l'entraînement, celui-ci est relativement faible car l'agent entre rapidement en collision avec les obstacles. Puis, après avoir appris les premiers comportements d'évitement, le temps moyen augmente : l'agent évite davantage les collisions mais peine à rejoindre efficacement la cible, ce qui conduit à davantage de timeout. Enfin, à mesure que la politique s'améliore, le temps moyen décroît puis se stabilise, traduisant la capacité de l'agent à atteindre son objectif par des trajectoires plus efficaces. \textbf{La récompense totale converge progressivement au fil de l'entraînement et la récompense de progression devient le terme dominant, l'agent privilégiant une progression régulière vers sa cible.}

\vspace{10pt}

\noindent \textbf{À partir des dernières expériences du curriculum, qui introduisent un aléa dans la position ou la taille des obstacles, les performances deviennent plus variables.} Les taux de collision et de timeout augmentent légèrement en raison de la diversité des configurations rencontrées, mais le taux de réussite reste globalement élevé. Ce comportement est toutefois attendu, ces scénarios présentant un niveau de difficulté supérieur.


%%%%%%%%%%%%%%%%%%%%%%%%%%%%%%%%%%%%%%%%%%%%%%%%%%%%%%%%%%%%%%%%%%%%%%%%%%%%%%%%%%%%%%%
%%%%%%%%%%%%%%%%%%%%%%%%%%%%%%%%%%%%%%%%%%%%%%%%%%%%%%%%%%%%%%%%%%%%%%%%%%%%%%%%%%%%%%%
%%%%%%%%%%%%%%%%%%%%%%%%%%%%%%%%%%%%%%%%%%%%%%%%%%%%%%%%%%%%%%%%%%%%%%%%%%%%%%%%%%%%%%%


\subsection{Analyse des scénarios d'entraînement}

\noindent Pour chaque expérience, les principaux indicateurs de performance sont relevés sur la politique finale obtenue après l'entraînement. Ces résultats permettent d'évaluer la capacité de l'agent à résoudre le scénario considéré avant d'introduire un niveau de difficulté supplémentaire.

\begin{longtable}{crrrrr}

    \caption{Performances obtenues à l'issue de l'entraînement} \label{tab:training_perf}                                              \\
    \toprule
    \textit{Expérience} & \textit{Steps}   & \textit{Succès (\%)} & \textit{Collision (\%)} & \textit{Temps (steps)} & \textit{Return} \\
    \midrule
    \endfirsthead
    \toprule
    \textit{Expérience} & \textit{Épisode} & \textit{Succès (\%)} & \textit{Collision (\%)} & \textit{Temps (steps)} & \textit{Return} \\
    \midrule
    \endhead
    \midrule
    \endfoot
    \bottomrule
    \endlastfoot

    Exp. 1              & 50 000           & 98                   & 1                       & 30                     & 1 093           \\
    \midrule
    Exp. 2              & 150 000          & 84                   & 9                       & 38                     & 551             \\
    \midrule
    Exp. 3              & 200 000          & 86                   & 4                       & 59                     & 990             \\
    \midrule
    Exp. 4              & 250 000          & 70                   & 10                      & 86                     & 873             \\
    \midrule
    Exp. 5              & 300 000          & 56                   & 32                      & 63                     & 708             \\
    \midrule
    Exp. 6              & 300 000          & 88                   & 12                      & 22                     & 468             \\
    \midrule
    Exp. 7              & 300 000          & 52                   & 46                      & 32                     & 635             \\
    \midrule
    Exp. 8              & 600 000          & 96                   & 4                       & 36                     & 1 059           \\
\end{longtable}
\vspace{-10pt}
{
    \footnotesize
    \color{darkgray}
    \noindent Lecture du tableau : Chaque ligne correspond à une expérience du curriculum learning. La colonne \textit{Steps} indique le nombre d'itérations nécessaires pour atteindre la politique finale présentée. Les performances sont ensuite calculées en évaluant cette politique sur 500 épisodes du scénario correspondant. Le taux de succès, le taux de collision, le temps moyen de trajet et le return moyen permettent ainsi de caractériser les performances finales obtenues.
}

\vspace{10pt}

\noindent \textbf{Les résultats de la Table \ref{tab:training_perf} montrent que les performances finales restent élevées sur l'ensemble des scénarios malgré une augmentation progressive de la difficulté.} Comme attendu, les expériences les plus complexes conduisent à une diminution du taux de succès et à une augmentation du taux de collision. Cette évolution s'explique par l'introduction progressive d'obstacles plus nombreux, plus volumineux ou générés aléatoirement, qui rendent la navigation plus difficile.

\begin{longtable}{crrrrrrrrr}

    \caption{Évaluation des politiques apprises à chaque étape du curriculum learning} \label{tab:curriculum_perf} \\
    \toprule
    \textit{Politique} & Exp. 1 & Exp. 2 & Exp. 3 & Exp. 4 & Exp. 5 & Exp. 6 & Exp. 7 & Exp. 8                     \\
    \midrule
    \endfirsthead
    \toprule
    \textit{Politique} & Exp. 1 & Exp. 2 & Exp. 3 & Exp. 4 & Exp. 5 & Exp. 6 & Exp. 7 & Exp. 8                     \\
    \midrule
    \endhead
    \midrule
    \endfoot
    \bottomrule
    \endlastfoot

    $\pi_1$            & 98     & 36     & 9      & 6      & 11     & 58     & 23     & 8                          \\
    \midrule
    $\pi_2$            & 44     & 80     & 75     & 61     & 25     & 43     & 13     & 0                          \\
    \midrule
    $\pi_3$            & 88     & 91     & 85     & 83     & 46     & 52     & 26     & 0                          \\
    \midrule
    $\pi_4$            & 80     & 75     & 72     & 68     & 32     & 52     & 18     & 2                          \\
    \midrule
    $\pi_5$            & 96     & 85     & 84     & 71     & 54     & 63     & 32     & 3                          \\
    \midrule
    $\pi_6$            & 96     & 87     & 84     & 82     & 53     & 88     & 50     & 21                         \\
    \midrule
    $\pi_7$            & 96     & 87     & 84     & 82     & 53     & 88     & 50     & 21                         \\
    \midrule
    $\pi_8$            & 100    & 95     & 84     & 80     & 62     & 91     & 59     & 99                         \\
\end{longtable}
\vspace{-10pt}
{
    \footnotesize
    \color{darkgray}
    \noindent Lecture du tableau : La politique $\pi_3$ correspond à la politique apprise après l'expérience 3. Ainsi, une valeur de 91 à l'intersection de la ligne $\pi_3$ et de la colonne Exp. 2 signifie que cette politique obtient un taux de réussite de 91\% lorsqu'elle est évaluée sur le scénario de l'expérience 2, pour 100 épisodes.
}

\vspace{10pt}

\noindent \textbf{La Table \ref{tab:curriculum_perf} montre une amélioration progressive des performances des politiques intermédiaires sur les scénarios les plus complexes.} On observe également que les performances restent élevées sur les scénarios déjà maîtrisés, ce qui montre que les compétences acquises au cours des premières étapes sont globalement conservées lors de l'apprentissage des tâches suivantes. Ces résultats confirment l'intérêt du curriculum learning.


%%%%%%%%%%%%%%%%%%%%%%%%%%%%%%%%%%%%%%%%%%%%%%%%%%%%%%%%%%%%%%%%%%%%%%%%%%%%%%%%%%%%%%%
%%%%%%%%%%%%%%%%%%%%%%%%%%%%%%%%%%%%%%%%%%%%%%%%%%%%%%%%%%%%%%%%%%%%%%%%%%%%%%%%%%%%%%%
%%%%%%%%%%%%%%%%%%%%%%%%%%%%%%%%%%%%%%%%%%%%%%%%%%%%%%%%%%%%%%%%%%%%%%%%%%%%%%%%%%%%%%%


\subsection{Évaluation de la politique}

\noindent La politique $\pi_8$ obtenue à l'issue du curriculum learning est ensuite évaluée sur différents scénarios de test afin d'analyser sa capacité de généralisation. Ces scénarios couvrent plusieurs niveaux de difficulté et permettent d'étudier le comportement de l'agent dans des environnements proches ou plus éloignés de ceux rencontrés pendant l'entraînement.

\begin{longtable}{crrrr}

    \caption{Performances obtenues à l'issue de l'évaluation} \label{tab:eval_perf}                                 \\
    \toprule
    \textit{Évaluation} & \textit{Succès (\%)} & \textit{Collision (\%)} & \textit{Temps (steps)} & \textit{Return} \\
    \midrule
    \endfirsthead
    \toprule
    \textit{Évaluation} & \textit{Succès (\%)} & \textit{Collision (\%)} & \textit{Temps (steps)} & \textit{Return} \\
    \midrule
    \endhead
    \midrule
    \endfoot
    \bottomrule
    \endlastfoot

    Éval. 1             & 77                   & 23                      & 29                     & 884             \\
    \midrule
    Éval. 2             & 56                   & 44                      & 22                     & 663             \\
    \midrule
    Éval. 3             & 96                   & 4                       & 36                     & 1 059           \\
    \midrule
    Éval. 4             & 58                   & 42                      & 29                     & 723             \\
    \midrule
    Éval. 5             & 76                   & 24                      & 10                     & 407             \\
\end{longtable}
\vspace{-10pt}
{
    \footnotesize
    \color{darkgray}
    \noindent Lecture du tableau : Chaque ligne correspond à un scénario d'évaluation. Les performances sont obtenues en évaluant la politique finale sur 500 épisodes du scénario considéré. Le taux de succès, le taux de collision, le temps moyen de trajet et le return moyen permettent de caractériser les performances de la politique sur chacun des environnements de test.
}

\vspace{10pt}

\noindent \textbf{Les résultats de la Table \ref{tab:eval_perf} montrent que la politique finale conserve de bonnes performances sur les environnements statiques, y compris lorsque les configurations diffèrent de celles rencontrées durant l'entraînement.} Les scénarios intégrant des obstacles dynamiques conduisent en revanche à une diminution du taux de succès et à une augmentation du taux de collision. Ce comportement est attendu, puisque ces interactions n'ont pas été prises en compte lors de l'entraînement. Les performances obtenues montrent néanmoins que la politique apprise constitue une base solide, qui pourra être enrichie par l'introduction progressive d'obstacles dynamiques au cours des prochaines expériences.
